\chapter{绪论：从设计模型到制造过程}
\label{ch:introduction}

\section{CAD、CAE、CAM各自解决什么问题}
\dfn{CAX与产品生命周期}{CAX是计算机辅助技术的统称。\textbf{CAD}（Computer Aided Design）辅助创建、修改、分析和优化设计，几何造型是其基础；\textbf{CAE}（Computer Aided Engineering）对产品进行数值分析和仿真，为设计评价与优化提供依据；\textbf{CAM}（Computer Aided Manufacturing）辅助制造过程的规划、管理和控制。广义CAD有时包括CAD与CAE，答题时应先说明采用的范围。}

三者共享产品数据，但关注的问题不同。例如设计一个承载支架：CAD确定孔径、壁厚和装配位置；CAE计算给定材料、约束和载荷下的应力与变形；CAM确定加工工艺、刀具轨迹和设备控制程序。\keypoint{几何模型是连接设计、分析与制造的数据基础，但单有外形不足以完成分析或制造。}

\begin{center}
\renewcommand{\arraystretch}{1.3}
\begin{tabularx}{\linewidth}{lXX}
\toprule
技术 & 主要输入与任务 & 典型输出 \\
\midrule
CAD & 设计要求、尺寸与约束；几何造型、工程图、装配、干涉检查 & 产品几何模型、装配模型、设计文档 \\
CAE & 几何、材料、网格、边界条件、载荷；应力、传热、流动、运动等分析 & 响应场、性能指标、优化建议 \\
CAM & 产品与工艺信息、设备资源；工艺规划、NC编程、机器人或CMM编程 & 工艺文件、刀位数据、设备控制程序 \\
\bottomrule
\end{tabularx}
\end{center}
\sourcepage{第1章}{3--4、19--21}

\reviewfig[.88]{ch01-lifecycle.pdf}{产品周期中的设计、分析和制造环节（摘自第1章讲义，PDF第4页）。}{fig:lifecycle}

\subsection{产品开发是一条反馈链}
设计需求经过信息收集与可行性研究，转化为概念设计、分析模型、分析与优化、设计评估及设计文档；随后进入工艺规划、生产准备、生产和质量控制。分析结果可能要求重新设计，加工能力也可能反过来约束结构尺寸，因此流程不应理解为只向前传递的单向流水线。

\ex{为什么相同的CAD模型会得到不同CAE结果？}{几何相同不代表分析问题相同。材料参数、网格密度、载荷大小与作用位置、约束方式不同，都会改变数值模型及其结果。复习时应把\textbf{几何模型}与\textbf{包含物理假设的分析模型}区分开。}

\section{CAD/CAM系统的组成与集成}
\subsection{硬件与软件}
硬件由主机、输入设备、显示设备、外存储设备和输出设备组成。主机承担计算与逻辑推理；键盘、鼠标、扫描仪等提供输入；显示器支持交互；CAD输出常为图纸或文档，CAM输出则进一步面向机床及其控制设备。

软件可分为三层：\textbf{系统软件}管理和维护计算机运行环境；\textbf{应用软件}完成专业设计、分析、制造任务；\textbf{二次开发软件}通过接口、脚本或扩展程序适配具体业务。三层是功能分工，不能把所有运行在计算机上的程序都称为CAD应用软件。
\sourcepage{第1章}{22--25}

\subsection{集成的关键是数据及其关联}
CAD/CAM可以减少重复建模与重复输入，缩短周期，提高设计一致性和产品标准化程度。其成立的前提是数据能被后续环节理解：例如孔的几何形状、尺寸公差、材料及加工要求必须有效传递，设计修改还应触发相关分析、工艺或程序的更新。

\dfn{PDM与并行工程}{\textbf{产品数据管理PDM}（Product Data Management）组织与管理产品开发数据及其相关过程，使CAD/CAM与生产、质量、资源计划等系统能够围绕受控数据协作。\textbf{并行工程}强调设计、工艺、制造等工作在协调下并行展开，尽早发现下游约束，减少后期返工。}

工业互联网支持跨地域的数据共享及制造资源协同；VR/AR用于设计展示、虚拟装配与信息叠加；人工智能可辅助生成设计、参数优化、建模与验证。这些应用均依赖可靠的数据、明确的目标与约束。
\sourcepage{第1章}{29--32}

\section{技术演进与软件功能识别}
讲义将发展分为技术准备、基础形成、全面发展与集成应用四个阶段。复习重点是\textbf{能力如何增长}：由交互绘图和数控加工起步，发展为二维与线框设计，再到实体造型、图形化编程和加工仿真，最终形成CAD/CAM集成及基于模型的数据贯通。MBD将几何与产品制造信息组织在三维模型中，第\ref{ch:cad-modeling}章会进一步讨论。

常见软件应按功能理解：二维制图与文档、参数化实体与装配、自由曲面、工程分析、数控加工等。选择工具时要看所需模型、数据交换能力与具体任务；软件名称本身不等于技术原理。智能化、集成化、虚拟化、网络化、绿色化和三维化是讲义归纳的发展方向。
\sourcepage{第1章}{27--37}

\subsection*{复习自测}
\begin{enumerate}
\item CAD、CAE、CAM的核心区别是什么？\textit{答：设计表达、性能分析、制造规划与控制。}
\item 为什么完成三维建模后还不能直接开展有限元计算？\textit{答：还需定义材料、网格、边界条件与载荷等。}
\item CAM是否仅指数控编程？\textit{答：否，还包括制造过程的规划、管理与控制，以及机器人、测量设备等编程。}
\item 集成系统为什么强调同一产品数据及修改关联？\textit{答：避免重复输入和版本不一致，使下游工作能跟随设计变化。}
\end{enumerate}
